\chapter{Methodology}\label{ch:methodology}
\section{Research Questions and Objectives}
The experiment is a controlled development-set ablation. It asks:
\begin{enumerate}
\item Does a cross-encoder improve the ranked evidence and downstream context relative to hybrid BM25 plus dense retrieval?
\item Does bounded traversal of source-backed document links add useful evidence beyond the same hybrid baseline?
\item Does combining graph retrieval and reranking improve quality enough to justify its observed inference latency?
\item When the context gate rejects evidence, does bounded retrieval feedback recover additional completed answers?
\end{enumerate}
These questions concern the implemented system and selected corpus. They are not hypotheses about all Vietnamese legal tasks. The study records quality at three stages and timing separately, rather than choosing a winner from one unqualified composite score.

\section{Data Collection and Preprocessing}
\subsection{Source Selection and Representation}
The prepared corpus resolves the source dataset to commit
\begin{quote}\small\path{a1b6860c821b7142ffe2385178af4e4bb6f86aff}.\end{quote}
Its local corpus identity is \code{4d70694000eb63e3}. Selection retains 7,425 article records from 647 originating documents in three topics. Metadata includes source dataset and revision, document identifiers and numbers, source URLs, and available article/clause/point locations. Unknown legal status is retained as unknown rather than interpreted as currently effective.

The hierarchical chunker uses a 512-token target and 64-token overlap for fallback token windows. It preserves article roots, recognizes numbered clauses, and splits long clauses into points when numbering is unambiguous. Ambiguous or oversized segments fall back to token windows. Heading context accompanies leaves. Parents are stored for later context expansion but are not embedded. This produces 14,708 embedded leaves and 14,333 parent chunks.

\begin{figure}[H]\centering
\begin{tikzpicture}[node distance=7mm]
\node[box,text width=4.2cm] (r) {Source article + document metadata\\unembedded parent};
\node[box,text width=3.5cm,below left=9mm and 5mm of r] (c1) {Complete clause parent\\conditions retained};
\node[box,text width=3.5cm,below right=9mm and 5mm of r] (c2) {Short provision leaf\\embedded directly};
\node[box,text width=3.5cm,below=of c1] (l) {Point / token-window leaves\\embedded with parent links};
\draw[arrow] (r)--(c1);\draw[arrow] (r)--(c2);\draw[arrow] (c1)--(l);
\node[below=8mm of l,font=\small,align=left,text width=11.5cm,anchor=north west,xshift=-1mm] {Retrieval ranks leaves. Context packing can expand a retained leaf to a complete parent when the token budget permits. Diagram shows representation rules, not a measured example document.};
\end{tikzpicture}\vspace{10mm}
\caption{Hierarchical representation separates retrieval granularity from answering context.}\label{fig:hierarchy}
\end{figure}

\subsection{Question Set and Evidence Labels}
The normalized file \code{three_topics_pilot.development.v1.jsonl} contains 45 questions: 15 per topic. Question types are shown in Figure~\ref{fig:types}. Each record contains a stable identifier, Vietnamese question, reference answer, atomic reference claims, and source evidence with document identity, legal location, and an exact quotation. There are 233 reference claims and 115 evidence spans. All labeled spans resolve against the prepared corpus.

\begin{figure}[H]\centering
\begin{tikzpicture}
\begin{axis}[xbar,width=.90\textwidth,height=5.0cm,xmin=0,xmax=30,xlabel={Number of questions},symbolic y coords={Multi-document,Comparison,Lookup,Multi-clause},ytick=data,nodes near coords,bar width=12pt,enlarge y limits=.2]
\addplot[fill=navy,draw=navy] coordinates {(3,Multi-document) (3,Comparison) (12,Lookup) (27,Multi-clause)};
\end{axis}
\end{tikzpicture}
\caption{Development-set composition: 45 questions, all labeled answerable; three topics contribute 15 questions each.}\label{fig:types}
\end{figure}

Twelve question wordings were clarified during development to align the requested scope with the recorded evidence. The original draft is preserved separately. The normalized file is therefore a development artifact, and the report does not treat the questions as a held-out test. The labels remain pending independent legal review.

The evidence resolver normalizes Unicode NFC and whitespace and matches quotations within the labeled document and provision location. It produces binary relevant leaf IDs. A leaf overlapping the exact labeled quote can be relevant, but alternative useful provisions are not exhaustively labeled. In this run, 39 questions have one relevant leaf and six have two. This narrow labeling scheme is essential to interpreting precision.

\section{Feature Selection and Engineering}
The project uses engineered representations and metadata rather than training tabular features. Lexical retrieval applies Unicode-aware, lowercase word tokenization and BM25Okapi with $k_1=1.2$ and $b=0.75$. Dense retrieval embeds question and leaf text into normalized 1,024-dimensional BGE-M3 vectors. Cosine search is served by pgvector HNSW. Exact document numbers and recognized legal locations are additional constraints, not embedding dimensions.

The original user question controls hard filters. LLM rewriting can improve retrieval wording, but cannot legitimately remove a document or condition from the request. Explicit document numbers are mapped to corpus vocabulary using deterministic normalization. Named article titles are used as hard constraints only when they can be resolved consistently for all referenced documents. Context prioritization can perform softer title matching and distribute candidates across requested documents, while leaving the recorded raw ranking intact for IR scoring.

For ranked lists $L_1,\ldots,L_m$, fusion assigns a candidate $d$ the score
\begin{equation}
\operatorname{RRF}(d)=\sum_{j:d\in L_j}\frac{1}{60+\operatorname{rank}_{L_j}(d)}.
\end{equation}
The optional reranker then scores query--candidate pairs. Graph candidates use the same metadata filters and enter fusion as an additional ranked list. The experiment fixes these shared rules so that the optional components can be compared without changing the source corpus or answer model.

\section{Model Training and Validation}
\subsection{Frozen Pretrained Inference}
No model is trained or fine-tuned in this study. The heading is retained from the course outline, but the relevant methodological choices are pretrained model selection, prompt configuration, and development validation. Table~\ref{tab:models} gives the measured-run roles. The online sufficiency evaluator and the offline quality judge have different responsibilities and must not be conflated.

\begin{table}[H]\centering\small
\caption{Models and budgets in the recorded development run.}\label{tab:models}
\begin{tabularx}{\textwidth}{P{3.5cm}YP{3cm}}\toprule
\textbf{Role} & \textbf{Model / backend} & \textbf{Key budget}\\\midrule
Dense embedding & BGE-M3 via Ollama & 1,024 dimensions\\
Query analysis & \code{qwen3.5:2b}, Ollama & 512 output tokens\\
Online context sufficiency & \code{qwen3.5:4b}, Ollama & 768 output tokens\\
Answer generation & \code{qwen3.5:4b}, Ollama & 2,048 output tokens\\
Citation support verification & \code{qwen3.5:2b}, Ollama & 512 output tokens\\
Optional reranking & mMARCO mMiniLMv2 L12 H384, ONNX int8 & 512 input tokens\\
Offline quality judging & \code{qwen3.5:2b}, Ollama & 2,048 output tokens\\\bottomrule
\end{tabularx}
\note{LLM context windows are 8,192 tokens; temperature is 0 and thinking is disabled. The offline judge differs from the generator model size, but shares the model family and the online verifier model. This is not independent human assessment.}
\end{table}

The reranker uses revision \code{1427fd652930e4ba29e8149678df786c240d8825}. The BGE-M3 tokenizer is pinned to revision \code{5617a9f61b028005a4858fdac845db406aefb181}. Model identities are captured in the run manifest. The configured backends and recorded digest are more informative than assuming the same mutable model tag always resolves to identical weights.

\subsection{Experimental Control}
\begin{table}[H]\centering\small
\caption{Controlled ablation matrix; every method includes BM25, dense retrieval and RRF.}\label{tab:matrix}
\begin{tabularx}{\textwidth}{lYccc}\toprule
\textbf{ID} & \textbf{Configuration} & \textbf{Graph} & \textbf{Rerank} & \textbf{Max. rounds}\\\midrule
B0 & Hybrid baseline & No & No & 1\\
B1 & Hybrid + reranker & No & Yes & 1\\
B2 & Hybrid + graph & Yes & No & 1\\
B3 & Hybrid + graph + reranker & Yes & Yes & 1\\
B4 & B3 + bounded refinement & Yes & Yes & 3\\\bottomrule
\end{tabularx}
\end{table}
All methods use 20 candidates per lexical/vector branch, a final ranked list of up to 20, and up to five packed context chunks. The context budget is 4,096 tokens with a 1,024-token prompt reserve. B4 preserves the configured methods and cutoffs across rounds; feedback refines the retrieval query rather than changing the ablation.

The run uses one repeat, seed 42, one excluded warm-up query per configuration, shuffled configuration order, and full query analysis for each timed execution. There are $45\times5=225$ expected executions. Warm-up reduces some startup effects but cannot remove local model-loading or scheduling variation. Neither repeated-run stability nor an independent test split is estimated by this pilot.

\section{Evaluation Metrics}
\subsection{Stage 1: Retrieval}
For question $q$, let $R_q$ be the unique labeled relevant leaf IDs and $z_i$ indicate whether the unique hit at rank $i$ belongs to $R_q$. For a cutoff $k$,
\begin{align}
P_q@k&=\frac{\sum_{i=1}^{k}z_i}{k},\\
R_q@k&=\frac{\sum_{i=1}^{k}z_i}{|R_q|},\\
RR_q@k&=\begin{cases}1/\min\{i\leq k:z_i=1\},&\text{if a relevant hit exists},\\0,&\text{otherwise},\end{cases}\\
AP_q@k&=\frac{\sum_{i=1}^{k}(P_q@i)z_i}{|R_q|},\\
DCG_q@k&=\sum_{i=1}^{k}\frac{z_i}{\log_2(i+1)},\\
nDCG_q@k&=\frac{DCG_q@k}{\sum_{i=1}^{\min(k,|R_q|)}1/\log_2(i+1)}.
\end{align}
Missing ranks contribute zero. The implementation deduplicates IDs while preserving order. MRR and MAP are macro means of $RR_q$ and $AP_q$, respectively; AP uses the total number of relevant IDs as its denominator, so missed labels remain misses. The report emphasizes $k=5$ and also reports recall at 20. It evaluates the final retained retrieval ranking for each question, not the context builder's reordered or expanded passages.

There are only 51 relevant question--leaf assignments across 45 questions. Even a perfect top-five ranking can therefore achieve at most
\begin{equation}
\overline{P@5}_{\max}=\frac{39\times1+6\times2}{45\times5}=0.2267.
\end{equation}
A precision value near 0.21 should be read against this label ceiling. It is neither a 21\% question-answer accuracy nor proof that every unlabeled retrieved passage is useless.

\subsection{Stage 2: Packed Context}
Let $v_i$ denote the offline judge's binary usefulness judgment for packed context chunk $i$, relative to the question and reference answer. Context Precision is rank-aware:
\begin{equation}
CP_q=\frac{\sum_{i=1}^{m}v_i\left(\sum_{j=1}^{i}v_j/i\right)}{\sum_{i=1}^{m}v_i},
\end{equation}
with zero when no context is judged useful. Its denominator counts useful \emph{retrieved} contexts. Consequently, a high score does not imply that all reference information was retrieved. It also differs from raw chunk-label $P@k$.

If the reference contains $s$ atomic claims and $u_j$ indicates whether claim $j$ is supported by the packed context, then Context Recall is $CR_q=\sum_j u_j/s$. These judgments operate on the actual packed text, including complete parent expansion when used. A correct leaf ID can be found while the semantic judge still considers the overall context incomplete.

\subsection{Stage 3: Answer Quality}
Faithfulness is the proportion of extracted answer claims judged supported by the supplied context. Factual precision is the fraction of response claims supported by the reference answer; factual recall is the fraction of reference claims covered by the response. Factual F1 is their harmonic mean, computed per question and then macro-averaged. It is not the harmonic mean of the two reported macro means.

Answer Relevancy averages cosine similarity between the original question embedding and three questions generated back from the answer, with noncommittal responses penalized by the scoring rule. It measures semantic alignment, not legal truth. All these measures depend on the project-specific extraction and entailment prompts. The additional \textit{correct-answer yield} requires completion, factual F1 equal to one, and faithfulness equal to one; insufficient-evidence outcomes contribute zero and invalid judge outputs are excluded. It is a strict judge-dependent diagnostic, not general accuracy.

\subsection{Inference Latency and Statistical Comparison}
Latency is wall-clock full-pipeline time per question, including query analysis, retrieval rounds, context evaluation, generation and verification when those stages execute. Index preparation, initial pipeline/model/tokenizer construction, excluded warm-up, and subsequent offline judging are outside this measurement. Backend request spans and model-call spans overlap with stage spans, so they must not be summed as independent costs. Results report means over all 45 outcomes and completed-only means, medians and 95th percentiles separately.

The benchmark pairs results by question, uses 2,000 question-level bootstrap resamples for pointwise 95\% confidence intervals, and applies Monte Carlo sign-flip tests with Holm adjustment across 234 reported metric/comparison tests. This broad exploratory family includes several correlated metrics. Answer metrics use only valid paired questions. A non-significant result does not establish equivalence; the small sample and one repeat limit power. Pointwise intervals are not family-adjusted evidence of superiority.

\clearpage
\section{Implementation Plan}
The benchmark is isolated from production orchestration. The validation step checks dataset structure and, when provided a corpus, resolves exact evidence. The run step records inference artifacts without invoking the offline quality judge. The score step adds deterministic retrieval metrics and optional model judgments. The report step aggregates stored scores and exports statistics. This separation permits inspection and rescoring without repeating costly inference or embedding preparation.

\begin{figure}[H]\centering
\begin{tikzpicture}[node distance=5mm]
\node[box,text width=2.6cm] (v) {Validate\\questions + evidence};
\node[box,text width=2.6cm,right=5mm of v] (r) {Run\\samples + timings};
\node[box,text width=2.6cm,right=5mm of r] (s) {Score\\qrels + judgments};
\node[box,text width=2.6cm,right=5mm of s] (p) {Report\\means + paired tests};
\draw[arrow] (v)--(r);\draw[arrow] (r)--(s);\draw[arrow] (s)--(p);
\end{tikzpicture}
\caption{Independent benchmark workflow. Offline judging does not alter the stored inference answer.}\label{fig:benchmark}
\end{figure}

\section{Ethical Considerations}
The system presents information from supplied sources and is not a substitute for professional assessment of a specific dispute. It should not silently transform missing evidence into a confident legal conclusion. Source provenance and an explicit fallback help readers inspect the basis of an answer, but neither removes model error nor resolves the effective legal status of a source.

The evaluation set contains public legal-text questions rather than personal case files. Future deployment with user cases would need explicit handling of private information, retention, and access control. For the present report, the main ethical obligations are accurate attribution, transparent uncertainty, and faithful reporting of failures and judge limitations. No unsupported current-validity conclusion is inferred from an amendment edge or a completed pipeline status.
